\documentclass{article}
\usepackage{amsmath, amssymb, amsthm}

\setlength{\parindent}{0pt}

\newtheorem*{claim}{Claim}

\begin{document}

\begin{claim}
    Let $f : A \to B$ be a function where $A$ is nonempty. Then $f$ is injective 
    if and only if it has a left inverse.
\end{claim}

\begin{proof}
    Assume $f$ is injective. We construct a function $g : B \to A$ as follows.
    
    \medskip

    For each $b \in \operatorname{im}(f)$, there exists a unique $a \in A$ such that 
    $f(a) = b$. Then define $g(b) = a$.

    \medskip

    For each $b \notin \operatorname{im}(f)$, define $g(b) = a$ arbitrarily for some $a 
    \in A$.

    \medskip

    We now verify that $g$ is a left inverse of $f$. Let $x \in A$. Then $f(x) \in 
    \operatorname{im}(f)$, so by construction of $g$, $g(f(x)) = x$. Hence $g \circ f = 
    \operatorname{id}_A$.

    \medskip

    Assume a left inverse $g : B \to A$ exists such that $g \circ f = \operatorname{id}_A$. 
    Let $a_1, a_2 \in A$ be such that $f(a_1) = f(a_2)$. Applying $g$ to both sides 
    gives $g(f(a_1)) = g(f(a_2))$. Since $g \circ f = \operatorname{id}_A$, this gives $a_1 
    = a_2$. Therefore $f$ is injective.

    \medskip

    Hence $f$ is injective if and only if it has a left inverse.
\end{proof}

\end{document}
